\section{调查方案设计}

\subsection{调查对象及范围界定}

\subsubsection{调查对象}

本研究以居民对城市环境噪声污染的主观感知与治理诉求为核心研究内容，将调查对象确定为浙江省杭州、宁波、温州三市常住居民。

本文所指常住居民，为在对应城市连续居住6个月及以上、拥有固定住所、日常主要活动均在本市的人群，涵盖本地户籍人口与外来常住人口。该界定标准排除游客、短期出差人员等临时流动群体，保证调研数据可以真实反映长期身处本地声环境中的居民主观体验。

为提升样本代表性，本次抽样不对受访者的年龄、职业、居住区域进行硬性限定，尽可能覆盖各类人口特征群体。年龄划分五个区间：18岁及以下、19-30岁、31-45岁、46-60岁、61岁及以上；职业包含机关事业单位职工、企业从业者、自由职业者、学生、退休人员等九大类；居住区域分为主城区、城市郊区、县城、乡镇、农村，以此充分区分不同生活场景下居民的噪声暴露差异。

\subsubsection{调查范围}

杭州、宁波、温州三地。

范围选取原因如下：首先，在民意诉求层面，三地在“民呼我为”平台的噪声投诉量占全省比重较高，居民诉求集中，民意反馈具有较强代表性；其次，在监测数据支撑层面，结合全省环境噪声监测数据，杭甬温三市的声环境质量与达标情况更具典型性，与前期现状分析中的典型问题高度契合；最后，在样本典型性层面，作为浙江省经济与人口规模较大的城市，其噪声污染类型多元、矛盾突出，能够为本次研究提供丰富且具有现实参考价值的样本支撑。

\subsection{调查方法}

经过小组讨论与分析，我们最终决定从政策、数据、问卷与访谈三个方面系统开展研究。

\subsubsection{文献调查法}

通过查阅国家及浙江省噪声污染防治相关法律法规、政策文件、统计年鉴与环境监测报告，梳理噪声污染治理的政策导向与现实背景，同时参考国内外声环境感知与噪声治理领域的相关研究成果，为本次调查的整体设计、问卷编制与分析框架提供理论依据与数据支撑。

\subsubsection{问卷调查法}

以杭州、宁波、温州三市常住居民为调查对象，采用线上问卷形式，围绕噪声来源、暴露程度、主观感受、影响评价与治理诉求等维度设计结构化问卷，收集居民对声环境的感知数据。通过统计分析，探究不同人群对噪声污染的感知差异与核心诉求，为后续问题分析与对策建议提供实证依据。

\subsubsection{访谈调查法}

选取部分典型投诉居民与社区工作人员开展半结构化访谈，深入了解噪声污染对居民生活的实际影响、治理过程中的现实困境及公众对噪声治理的意见建议，弥补问卷数据的局限性，挖掘民意反馈背后的深层矛盾，为研究提供更具情境性与深度的质性资料。

\subsection{研究设计}

本次研究框架图如下：

\fig{图片8.png}{研究设计}

\subsubsection{问卷设计}

问卷设计应遵循目的性、逻辑性、明确性、可接受性和效率性等原则，以确保数据收集科学、有效且易于分析。

本文围绕居民对环境噪声的感知、噪声暴露现状、噪声治理评价以及相关认知建议进行设计，在设计上，我们注重问题的逻辑递进与作答便捷性，以确保收集到的数据科学规范、真实有效，为后续的统计分析与问题研判提供坚实支撑。问卷主体共包含五个核心模块：基本信息、噪声暴露与环境感知、噪声治理满意度、噪声治理认知以及噪声治理建议。

\begin{table}[H]
\centering
\caption{问卷具体内容}
\begin{tabular}{M{5cm}M{7cm}}
\bluetoprule
\rowcolor{tabhead}
\thead{模块} & \thead{题项} \\
\hline
\multirow{8}{*}{基本信息} & 性别 \\ \cline{2-2}
 & 年龄 \\ \cline{2-2}
 & 居住地 \\ \cline{2-2}
 & 居住区域类型 \\ \cline{2-2}
 & 职业类型 \\ \cline{2-2}
 & 文化程度 \\ \cline{2-2}
 & 是否为隔音墙 \\ \cline{2-2}
 & 居家办公情况 \\ \hline
\multirow{6}{*}{噪声暴露与环境感知} & 受噪声影响程度 \\ \cline{2-2}
 & 隔音水平评价 \\ \cline{2-2}
 & 各类噪声干扰程度 \\ \cline{2-2}
 & 离噪声源距离 \\ \cline{2-2}
 & 噪音影响时段 \\ \cline{2-2}
 & 噪声负面影响 \\ \hline
\multirow{3}{*}{噪声治理满意度} & 声环境质量满意度 \\ \cline{2-2}
 & 噪声治理满意度 \\ \cline{2-2}
 & 主要应对方式 \\ \hline
\multirow{4}{*}{噪声治理认知} & 投诉渠道认知 \\ \cline{2-2}
 & 国家政策认知 \\ \cline{2-2}
 & 治理必要性认知 \\ \cline{2-2}
 & 责任主体认知 \\ \hline
\multirow{2}{*}{噪声治理建议} & 举措优先度 \\ \cline{2-2}
 & 具体意见 \\
\bluebottomrule
\end{tabular}
\end{table}

\subsubsection{访谈设计}

\textbf{（一）访谈对象}

本次访谈聚焦两类典型噪声易感群体，选取杭州、宁波、温州三地主城区老旧小区独居老人、郊区工厂周边一线职工作为访谈对象。

老旧小区独居老人与工厂一线职工，分别代表城市生活区与工业区的典型噪声敏感人群。前者居住环境隔音薄弱、对噪声高度敏感，后者长期暴露于工业、交通及施工噪声中，两类群体均能精准反映杭甬温三地噪声污染的民生痛点，具备高度代表性。

\textbf{（二）访谈内容}

对象一为主城区老旧小区独居老人，提纲如下：

\begin{enumerate}[label=（\arabic*）,leftmargin=2.5em,itemsep=0.5ex]
\item 您在此处居住多久？房屋墙体隔音效果如何？
\item 日常生活中受到哪些环境噪声影响，影响程度如何？
\item 邻居生活噪声、广场舞、犬吠声等是否对您造成干扰？
\item 您居住区域整体隔音水平如何？哪个时段噪声影响最大？
\item 噪声是否影响您的睡眠、精神状态，有无耳鸣、烦躁等不适？
\item 您对居住地声环境质量、本地噪声治理工作是否满意？
\item 您是否了解12345、浙里办等噪声投诉渠道，是否听说过《中华人民共和国噪声污染防治法》？
\item 您认为是否需要强化浙江省噪声治理？噪声治理的责任主体是谁？
\item 针对改善噪声污染，您有哪些具体意见与建议？
\end{enumerate}

对象二为郊区工厂周边一线职工，提纲如下：

\begin{enumerate}[label=（\arabic*）,leftmargin=2.5em,itemsep=0.5ex]
\item 您在此职工宿舍居住多久？墙体隔音条件如何？是否需要居家办公？
\item 日常生活中受噪声影响程度如何，哪种噪声干扰最为严重？
\item 装修施工、商业经营、公共活动类噪声是否对您产生干扰？区域隔音水平怎么样？
\item 距离噪声源远近如何？哪个时段噪声影响最大？
\item 噪声是否造成睡眠障碍、精神疲惫、注意力下降等问题？
\item 长期噪声环境是否引发耳鸣、心悸、头晕等身体不适？
\item 遭遇噪声扰民时，您及同事主要采取哪些应对方式？
\item 您是否了解12345、浙里办等投诉渠道，是否知晓《中华人民共和国噪声污染防治法》？
\item 您认为是否需要强化浙江省噪声治理？责任主体包含哪些？
\item 针对工厂周边噪声污染治理，您有哪些具体的意见与建议？
\end{enumerate}

\subsection{样本容量确定}

样本容量的科学确定是保障问卷调查数据可靠性、统计分析有效性的核心前提，直接影响后续实证研究结果稳定性与准确性。以下是确定样本容量的方法。

\subsubsection{问卷预发放}

为确保正式问卷的可靠性，本研究先抽取杭甬温居民通过预调查的方式先发放了108份电子问卷，以作答时间少于60秒作为无效问卷的判定标准，经初步筛选后得到有效问卷99份，则问卷有效率约为92\%。该结果为后续正式调查的样本量修正提供了现实依据，也验证了问卷整体作答体验良好，题项设置基本合理。具体情况如下：

\begin{table}[H]
\centering
\caption{问卷预调研回收情况}
\label{tab:presurvey}
\rowcolors{2}{tablight}{white}
\begin{tabular}{M{3.5cm}M{3.5cm}M{3.5cm}M{3.5cm}}
\bluetoprule
\rowcolor{tabhead}
\thead{入样城市} & \thead{实发问卷} & \thead{有效问卷} & \thead{有效率} \\
\hline
杭州市 & 32 & 30 & 93.8\% \\
宁波市 & 46 & 42 & 91.3\% \\
温州市 & 30 & 27 & 90.0\% \\
合计 & 108 & 99 & 91.7\% \\
\bluebottomrule
\end{tabular}
\end{table}

\subsubsection{计算样本容量}

本研究采用简单随机抽样下的总体比例估计法确定最小样本量，计算公式为：

\begin{equation}
n_0=\frac{z_{\alpha/2}^{2}\cdot p(1-p)}{e^{2}}
\end{equation}

其中 $z_{\alpha/2}$ 为95\%置信水平下对应的标准正态分布分位数，取值1.96，$p$ 表示总体比例最保守情况的估计值，取值0.5，$e$ 为允许抽样的边际误差，为了提高结果的精度，我们设定为0.05。代入公式我们得到初始样本量：

\begin{equation}
n_1=\frac{1.96^{2}\times0.5\times0.5}{0.05^{2}}\approx384
\end{equation}

本研究的调查对象为杭州、宁波、温州三市常住居民，根据最新人口统计数据，三市常住居民总数约为 $N=3230.9$ 万，由于采用有限总体抽样，需对初始样本量进行有限总体修正，修正公式为：

\begin{equation}
n_2=n_1\times\frac{N}{N+n_1-1}\approx384
\end{equation}

由于三市常住居民总数 $N$ 远大于初始样本量 $n_1$，因此修正后仍取原来的数目。本研究采用分区域随机抽样，问卷设计效益（deff）设为1.1，继续修正为：

\begin{equation}
n_3=n_2\times\mathrm{deff}\approx422.4
\end{equation}

基于预调查得到的问卷有效率为 $r=92\%$，对样本量进行最终修正，得到正式调查所需发放的问卷总数，公式为：

\begin{equation}
n_4=\frac{n_3}{r}=\frac{422.4}{0.92}\approx459
\end{equation}

为了进一步保证数据有效性，以及考虑到现实需求，我们最终决定发放问卷480份，具体各市问卷情况如表\ref{tab:presurvey}所示。

\subsection{问卷回收情况}

问卷发放周期为2026年7月9日至2026年7月25日，共发放问卷520份，回收问卷493份，问卷回收率为94.81\%。对回收问卷进行有效性筛选，剔除作答时长不足100秒、所有题项选择同一答案、量表题项存在明显作答矛盾的无效问卷后，最终得到有效问卷467份（具体如表\ref{tab:recovery}）。

有效样本中，男性260人，女性207人；年龄分布以18-35岁群体为主，占比61.24\%；文化程度以本科及以上学历为主，占比73.02\%；样本在居住城市、居住区域类型、职业状态等维度分布均衡，具有良好的代表性，可用于后续统计分析。

\begin{table}[H]
\centering
\caption{问卷最终回收情况}
\label{tab:recovery}
\rowcolors{2}{tablight}{white}
\begin{tabular}{M{3.5cm}M{3.5cm}M{3.5cm}M{3.5cm}}
\bluetoprule
\rowcolor{tabhead}
\thead{入样城市} & \thead{回收问卷} & \thead{有效问卷} & \thead{有效率} \\
\hline
杭州市 & 142 & 132 & 93.0\% \\
宁波市 & 203 & 192 & 94.6\% \\
温州市 & 148 & 143 & 96.6\% \\
合计 & 493 & 467 & 94.7\% \\
\bluebottomrule
\end{tabular}
\end{table}

\subsection{问卷误差分析}

本次针对杭甬温三地居民开展环境噪声感知问卷调查，在数据收集、整理过程中不可避免存在误差，为保障研究结果的准确性与可靠性，对两类误差来源进行系统分析，并制定对应控制措施。

\subsubsection{抽样误差}

抽样误差是因采用抽样调查而非全面普查，样本结构与总体真实情况存在偏差所产生的随机误差，属于不可避免的统计性误差。

本次调查通过线上发放电子问卷开展，选取杭州、宁波、温州各类典型居住区域的不同人群作为调查样本。受样本覆盖范围限制，样本在年龄、职业、居住区域、居住年限等维度的分布，无法完全匹配区域全体居民的真实结构；同时受样本量限制，部分小众群体的噪声感知特征未能得到充分体现，进而导致样本统计结果与总体真实水平存在小幅偏差。这类误差可以通过合理扩大样本量、优化抽样结构降低，但无法完全消除。

\subsubsection{非抽样误差}

非抽样误差是由人为因素引发的系统性偏差，对本次调研数据质量的影响较为突出。

其来源主要集中在实际作答与数据处理两个阶段：在实际作答环节，部分受访者存在敷衍心态，表现为填写时长过短、答案逻辑矛盾或量表选项单一重复等规律性作答行为，导致问卷有效性不足；同时，由于部分受访者对噪声相关概念认知模糊，常依靠主观感受随意作答，进一步造成了数据的客观失真。而在数据回收、筛选与录入阶段，无效问卷的误留存以及人工录入过程中的操作失误，也共同构成了本次调研的非抽样误差来源。

\subsubsection{误差控制}

为有效控制两类误差对研究结论的干扰，本研究从抽样、作答到数据处理的全流程实施了较为系统的误差控制措施。

\textbf{（1）在调查设计层面}，参考同类环境感知问卷的成熟量表，结合预调查反馈优化题项表述，降低理解歧义，采用分层抽样设计以降低抽样误差，合理控制问卷长度，避免受访者因疲劳产生系统性偏差。

\textbf{（2）在数据收集层面}，同步通过线上、线下渠道采集样本，覆盖不同数字化能力群体，在问卷中穿插陷门题与方向反转题，识别并剔除随意作答的无效问卷，实时监控各城市样本采集进度，动态调整数据收集策略。

\textbf{（3）在数据后处理层面}，比对样本分布与总体人口学特征，按需采用事后分层加权修正样本偏差，运用多重插补法处理可控缺失变量，对存在系统性缺失的变量进行审慎处置，全部统计分析均依托清洗核验后的高质量数据开展。通过以上分层管控手段，本研究有效保障调查数据的信度与效度，为后续描述性统计、组间差异检验、回归建模等定量分析工作奠定可靠数据基础。

